\newpage
\subsection {Diagrama de caso de uso}

Os requisitos foram agrupados, e na fase de análise foram definidos os casos de uso. A Figura~\ref{fig:diagramacasouso} representa o diagrama de caso de uso da aplicação.

\begin{figure}[!htbp]
	\centering
	\caption{Diagrama de caso de uso.}
	\fbox{\includegraphics[scale=.6]{imagens/exemploQuadro.png}}
    \legend{Fonte: Autoria própria (2026)}	
	\label{fig:diagramacasouso}
\end{figure}

\newpage
\subsection {Requisitos funcionais}

Na definição dos requisitos funcionais para o desenvolvimento da aplicação, foram agrupados conforme apresentado no Quadro~\ref{tab:casosreq}. 

\begin{quadro}[!htbp]
\centering
\caption{Casos de Uso x Requisitos funcionais}
    \begin{tabular}{ | m{1.7cm} | m{4cm}| m{9cm} | }
		\hline
		\cellcolor{gray!20}\textbf{No} & \cellcolor{gray!20}\textbf{Caso de Uso / Módulo} & \cellcolor{gray!20}\textbf{Descrição resumida} \\ \hline
		RF-001 & Autenticação Unificada (SSO) & Permitir cadastro e login por meio de credenciais próprias ou OAuth 2.0 (Google, Apple, entre outros), mantendo a sessão sincronizada entre Web e App. \\ \hline
		RF-002 & Gestão de Perfil e TDU & Permitir ao usuário gerenciar seus dados pessoais, histórico de buscas e refazer o quiz de Tipo de Uso (TDU) a qualquer momento. \\ \hline
      	RF-003 & Sincronização Cross-Device & Sincronizar dispositivos favoritados e histórico de comparações em tempo real entre todas as plataformas autenticadas. \\ \hline
      	RF-004 & Quiz de Perfilamento & Executar um fluxo interativo de perguntas (orçamento, finalidade principal e portabilidade) para classificar o usuário em um TDU específico. \\ \hline
      	RF-005 & Busca Semântica & Converter termos de linguagem natural (ex.: "celular com câmera boa e barato") em filtros técnicos estruturados no banco de dados. \\ \hline
      	RF-006 & Tags Baseadas em TDU & Exibir badges visuais nos produtos, priorizando atributos relevantes de acordo com o perfil do usuário (ex.: "Bateria Longa" para "Estudante"). \\ \hline
     	RF-007 & Comparação Multidimensional & Gerar uma matriz comparativa lado a lado de até 4 dispositivos, destacando visualmente os atributos superiores. \\ \hline
      	RF-008 & Dicionário Contextual & Exibir modais e tooltips explicativos para termos técnicos, contendo traduções leigas com apoio visual (ícones e descrições simples). \\ \hline
        RF-009 & Gestão de Entidades & Disponibilizar painel administrativo com operações CRUD para fabricantes, modelos, especificações técnicas e mídias. \\ \hline
        RF-010 & Gestão do Dicionário & Oferecer interface administrativa para criação e atualização das traduções leigas dos termos técnicos. \\ \hline
	\end{tabular}
    \label{tab:casosreq}
    \\ \vspace{0.2cm}
	Fonte: Autoria própria (2026)
\end{quadro}

\subsection {Requisitos de regras de negócio}
Os requisitos de regras de negócios são readequados para a aplicação e apresentados no Quadro~\ref{tab:reqregrasnegocio}.

\begin{quadro}[!htbp]
\centering
\caption{Requisitos de regras de negócio}
	\begin{tabular}{ | m{1.7cm} | m{3cm}| m{10cm} | }
		\hline
		\cellcolor{gray!20}\textbf{No} & \cellcolor{gray!20}\textbf{Área de negócio} & \cellcolor{gray!20}\textbf{Descrição resumida} \\ \hline
        RGN-01 & Perfilamento & O usuário pode refazer o quiz de Tipo de Uso (TDU), mas só pode ter um perfil principal ativo por vez.\\ \hline
        RGN-02 & Catálogo & Um produto só pode ser listado no aplicativo se contiver os dados técnicos essenciais exigidos para a busca semântica.\\ \hline
        RGN-03 & Afiliados & Todo redirecionamento para o varejo deve obrigatoriamente injetar as UTMs e o ID de afiliado da plataforma.\\ \hline
        RGN-04 & Autenticação & Contas originadas via OAuth 2.0 não possuem senha gerenciada no sistema, dependendo do token do provedor.\\ \hline
        RGN-05 & Backoffice & Apenas usuários com privilégios de Administrador podem criar, editar ou excluir termos no Dicionário Contextual.\\ \hline       
	\end{tabular}
	\\ \vspace{0.2cm}
	Fonte: Autoria própria (2026)
	\label{tab:reqregrasnegocio}
\end{quadro}


\subsection {Requisitos não funcionais}
Os requisitos não funcionais são listados no Quadro~\ref{tab:reqnaofunc1}.

\begin{quadro}[!htbp]
\centering
\caption{Requisitos não funcionais}
	\begin{tabular}{ | m{1.7cm} | m{4cm}| m{9cm} | }
		\hline
		\cellcolor{gray!20}\textbf{No} & \cellcolor{gray!20}\textbf{Quanto a(o):} & \cellcolor{gray!20}\textbf{Descrição resumida} \\ \hline
        RNF-01 & Segurança & O processamento de preferências, histórico e dados de sessão deve seguir as diretrizes da LGPD (Lei Geral de Proteção de Dados).\\ \hline
        RNF-02 & Sincronização & O histórico de buscas e dispositivos favoritados deve ser sincronizado em tempo real ou com baixíssima latência entre plataformas.\\ \hline
        RNF-03 & Disponibilidade & Os serviços de catálogo e redirecionamento de compras devem possuir alta disponibilidade (uptime de 99,9\%).\\ \hline
        RNF-04 & Compatibilidade & A plataforma deve ser responsiva na Web e possuir aplicativo perfeitamente funcional em ecossistemas Android e iOS.\\ \hline
        RNF-05 & Usabilidade & O dicionário de hardware e os modais explicativos devem seguir as normas de acessibilidade (WCAG) para contraste e leitores de tela.\\ \hline
	\end{tabular}
	\\ \vspace{0.2cm}
	Fonte: Autoria própria (2026)
	\label{tab:reqnaofunc1}
\end{quadro}
